\documentclass[9pt,aspectratio=169]{beamer}

\usepackage{graphicx}
% \usepackage[OT1, T1]{fontenc}
\usepackage{cmbright}
% \usepackage{sansmath}
\usepackage{amsmath}
\usepackage{mathrsfs}
\usepackage{caption}
\usepackage{bbm}
%\usepackage{enumitem}
\usepackage{natbib}
\usepackage{booktabs}
\definecolor{theme}{HTML}{2e24b1}
\usefonttheme{professionalfonts}
\DeclareSymbolFont{greekletters}{OML}{sansmath}{m}{it}
\DeclareMathSymbol{\beta}{\mathord}{greekletters}{"0C}
\usepackage{etoolbox}
%\DeclareMathSymbol{\delta}{\mathord}{greekletters}{"0E}
%\setbeamercolor{structure}{fg=theme}
% \usetheme{default}
%\usetheme{AnnArbor}
%\usetheme{Antibes}
%\usetheme{Bergen}
%\usetheme{Berkeley}
%\usetheme{Berlin}
%\usetheme{Boadilla}
%\usetheme{boxes}
%\usetheme{CambridgeUS}
%\usetheme{Copenhagen}
%\usetheme{Darmstadt}
\usetheme{Frankfurt}
%\usetheme{Goettingen}
%\usetheme{Hannover}
%\usetheme{Ilmenau}
%\usetheme{JuanLesPins}
%\usetheme{Luebeck}
%\usetheme{Madrid}
%\usetheme{Malmoe}
%\usetheme{Marburg}
%\usetheme{Montpellier}
%\usetheme{PaloAlto}
%\usetheme{Pittsburgh}
%\usetheme{Rochester}
%\usetheme{Singapore}
%\usetheme{Szeged}
%\usetheme{Warsaw}

% \usecolortheme{default}
% \usecolortheme{albatross}
% \usecolortheme{beaver}
% \usecolortheme{beetle}
% \usecolortheme{crane}
% \usecolortheme{dolphin}
% \usecolortheme{dove}
% \usecolortheme{fly}
% \usecolortheme{lily}
% \usecolortheme{orchid}
% \usecolortheme{rose}
% \usecolortheme{seagull}
\usecolortheme{seahorse}
% \usecolortheme{whale}
% \usecolortheme{wolverine}
\newcommand{\aA}{\mathscr A}
\newcommand{\cC}{\mathscr C}
\newcommand{\dD}{\mathscr D}
\newcommand{\bB}{\mathscr B}
\newcommand{\oO}{\mathcal O}
\newcommand{\gG}{\mathscr G}
\newcommand{\hH}{\mathcal H}
\newcommand{\kK}{\mathcal K}
\newcommand{\iI}{\mathcal I}
\newcommand{\mM}{\mathcal M}
\newcommand{\eE}{\mathcal E}
\newcommand{\fF}{\mathcal F}
\newcommand{\qQ}{\mathcal Q}
\newcommand{\tT}{\mathcal T}
\newcommand{\xX}{\mathcal X}
\newcommand{\yY}{\mathcal Y}
\newcommand{\wW}{\mathcal W}
\newcommand{\uU}{\mathcal U}
\newcommand{\lL}{\mathcal L}
\newcommand{\rR}{\mathcal R}
\newcommand{\zZ}{\mathcal Z}

\newcommand{\sS}{\mathscr S}

\newcommand{\pP}{\mathcal P}

\newcommand{\vV}{\mathcal V}

\newcommand{\BB}{\mathbbm B}
\newcommand{\RR}{\mathbbm R}
\newcommand{\CC}{\mathbbm C}
\newcommand{\QQ}{\mathbbm Q}
\newcommand{\NN}{\mathbbm N}
\newcommand{\GG}{\mathbbm G}
\newcommand{\UU}{\mathbbm U}
\newcommand{\ZZ}{\mathbbm Z}
\newcommand{\KK}{\mathbbm K}
\newcommand{\PP}{\mathbbm P}
\newcommand{\EE}{\mathbbm E}
\newcommand{\TT}{\mathbbm T}

\newcommand{\var}{\mathbbm V}
\renewcommand{\implies}{\Rightarrow}
\renewcommand{\iff}{\Leftrightarrow}

\newcommand{\II}{\mathbb I}
\newcommand{\WW}{\mathbb W}

\renewcommand{\phi}{\varphi}

\newcommand{\XX}{\mathsf X}
\newcommand{\YY}{\mathsf Y}

\newcommand{\bP}{\mathbf P}
\newcommand{\bQ}{\mathbf Q}
\newcommand{\bE}{\mathbf E}
\newcommand{\bM}{\mathbf M}
\newcommand{\bX}{\mathbf X}
\newcommand{\bY}{\mathbf Y}

\setbeamertemplate{theorems}[numbered]
% \newtheorem{lemma}[theorem]{Lemma}
\newtheorem{proposition}[theorem]{Proposition}

\renewcommand{\baselinestretch}{1.3}
\setlength{\parskip}{1.0ex plus0.5ex minus0.5ex}
\setbeamertemplate{caption}[numbered]
\setbeamertemplate{frametitle continuation}{}
\setbeamertemplate{headline}{}
\defbeamertemplate{footline}{myframe number}
{
  \hspace{1pt}
    \usebeamercolor[structure]{page number in head/foot}%
  \usebeamerfont{page number in head/foot}%
  \raisebox{1.7ex}[0pt][0pt]{% <--- change here
    \insertframenumber\,/\,\inserttotalframenumber\kern1em}%
}
\setbeamertemplate{footline}[myframe number]

\makeatletter
\patchcmd{\beamer@sectionintoc}
  {\vfill}
  {\vskip\itemsep}
  {}
  {}
\makeatother  

\AtBeginSection[]
{
  \begin{frame}
    \frametitle{Outline}
    \begin{minipage}{1.0\linewidth}
      \tableofcontents[currentsection,currentsubsection]
    \end{minipage}
  \end{frame}
}

\AtBeginSubsection[]
{
  \begin{frame}
    \frametitle{Outline}
    \begin{minipage}{1.0\linewidth}
      \tableofcontents[currentsection,currentsubsection]
    \end{minipage}
  \end{frame}
}


\title[Lecture 1: Basic Set Theory]{Lecture 1: Basic Set Theory}

\author{Junnan Zhang} % Your name
\institute[Paula and Gregory Chow Institute for Studies in Economics\\
Xiamen University] % Your institution as it will appear on the bottom of every slide, may be shorthand to save space
{
	Paula and Gregory Chow Institute for Studies in Economics \\ Xiamen University  \\ % Your institution for the title page
	\medskip
}
\date{Fall, 2026} % Date, can be changed to a custom date
\setbeamertemplate{itemize items}[default]
\setbeamertemplate{sidebar}[sidebar right]
% \definecolor{myco}{HTML}{a67678}
\definecolor{myco}{HTML}{8776a6}
\setbeamercolor{itemize item}{fg=myco} % all frames will have red bullets
\setbeamercolor{itemize subitem}{fg=myco} % all frames will have red bullets
\setbeamercolor{block title}{bg=myco}
\setbeamercolor{structure}{fg=myco}
% - Either use conference name or its abbreviation.
% - Not really informative to the audience, more for people (including
%   yourself) who are reading the slides online

\subject{Economics}
% This is only inserted into the PDF information catalog. Can be left
% out. 

% If you have a file called "university-logo-filename.xxx", where xxx
% is a graphic format that can be processed by latex or pdflatex,
% resp., then you can add a logo as follows:

% \pgfdeclareimage[height=0.5cm]{university-logo}{university-logo-filename}
% \logo{\pgfuseimage{university-logo}}

% Delete this, if you do not want the table of contents to pop up at
% the beginning of each subsection:
% \AtBeginSubsection[]
% {
%   \begin{frame}<beamer>{Outline}
%     \tableofcontents[currentsection,currentsubsection]
%   \end{frame}
% }

% Let's get started
\begin{document}

\begin{frame}
    \titlepage
\end{frame}

\begin{frame}{Summary}
    Topics to be covered today
    \begin{itemize}
        \item Sets
        \item Functions
        \item Real Numbers
    \end{itemize}
\end{frame}

\section{Sets}

\begin{frame}
    \frametitle{Introduction}
    \begin{itemize}
        \item Set theory is the language of all mathematics
        \item In economics, sets are everywhere: consumption sets, production
        sets, sets of players, sets of equilibria, etc.
    \end{itemize}
\end{frame}

\begin{frame}[allowframebreaks]{Sets}
    \begin{definition}
        \begin{itemize}
            \item A set is, informally speaking, a collection of objects, which
            we call \textbf{elements}.
            \item If $x$ is an element of the set $S$, we write $x \in S$. We
            also say $x$ \textbf{belongs to} $S$.

            If $x$ is not an element of $S$, we write $x \notin S$.
        \end{itemize}
    \end{definition}
    \framebreak
    We use curly brackets when describing a set. We can list all its
    elements:
    \begin{itemize}
        \item $\{1, 2, 3\}$
        \item $\big\{1, 2, \{1, 3\}\big\}$
    \end{itemize}
    Or we can specify the properties that its elements satisfy:
    \begin{itemize}
        \item $\{x: x \text{ is irrational}\}$: the set of all $x$ such
        that $x$ is irrational
        \item $\{n\in\NN: n^2 < 7\}$: the set of all natural numbers $n$
        such that $n^2 < 7$
    \end{itemize}
    \framebreak
    \begin{definition}
        \begin{itemize}
            \item We say two sets $A$ and $B$ are \textbf{equal} if they
            have the same elements, and we write $A = B$.
            \item An \textbf{empty set} is a set with no elements and it
            is denoted by $\emptyset$.
            \item If $\forall x\in A, x\in B$, we say $A$ is a
            \textbf{subset} of $B$ and we write $A \subset B$.
            \item The set of all subsets of $A$ is called the
            \textbf{power set} of $A$ and is denoted by $\pP(A)$.
        \end{itemize}
    \end{definition}
    We usually prove $A = B$ by proving $A \subset B$ and $B \subset A$.
\end{frame}

\begin{frame}{Sets}
    \framesubtitle{Examples}
    \begin{itemize}
        \item $\{\emptyset\}$ is not an empty set because $\emptyset \in
        \{\emptyset\}$
        \item $\emptyset \subset A$ for all sets $A$
        \item If $A = \{1, 2\}$, then $\pP(A) = \{\emptyset, \{1\}, \{2\},
        \{1, 2\}\}$
        \item $\emptyset \in \pP(A)$ for all sets $A$
    \end{itemize}
\end{frame}

\begin{frame}{Sets}
    \begin{proposition}
        Set inclusion is transitive: if $A \subset B$ and $B \subset C$, then
        $A \subset C$.
    \end{proposition}
    \emph{Proof:} 
\end{frame}

\begin{frame}{Union and Intersection of Sets}
    \begin{definition}
        \begin{itemize}
            \item The \textbf{union} of $A$ and $B$ is defined by
            \begin{equation*}
                A \cup B := \{x: x\in A \text{ or } x\in B\}
            \end{equation*}
            \item The \textbf{intersection} of $A$ and $B$ is defined by
            \begin{equation*}
                A \cap B := \{x: x\in A \text{ and } x\in B\}
            \end{equation*}
            \item Two sets $A$ and $B$ are \textbf{disjoint} if $A \cap B =
            \emptyset$ 
        \end{itemize}
    \end{definition}
\end{frame}

\begin{frame}{Union and Intersection of Sets}
    \framesubtitle{Venn Diagram}
\end{frame}


\begin{frame}[allowframebreaks]{Union and Intersection of Sets}
    \begin{definition}
        \begin{itemize}
            \item If $\fF$ is a family of sets, the union of all sets in
            $\fF$ is defined by
            \begin{equation*}
                \bigcup \fF := \{x: x\in A \text{ for at least one } A \in \fF\}
            \end{equation*}
            \item Similarly, the intersection of all sets in $\fF$ is defined by
            \begin{equation*}
                \bigcap \fF := \{x: x\in A \text{ for every } A \in \fF\}
            \end{equation*}
            \item If $\fF = \{A_\lambda: \lambda \in J\}$, then
            \begin{equation*}
                \bigcup \fF = \bigcup_{\lambda \in J} A_\lambda \qquad \bigcap \fF =
                \bigcap_{\lambda \in J} A_\lambda
            \end{equation*}
        \end{itemize}
    \end{definition}
    \framebreak
    \begin{proposition}
        Let $A, B, C$ and $B_\lambda$ (for $\lambda \in J$) be sets. Then
        \begin{align*}
            A \cup B = B \cup A &\quad A\cap B = B \cap A\\
            A \cup (B \cup C) = (A \cup B) \cup C &\quad
            A \cap (B \cap C) = (A \cap B) \cap C \\
            A \subset A \cup B &\quad A \cap B \subset A \\
            A \subset B \iff A\cup B = B &\quad
            A \subset B \iff A \cap B = A\\
            A \cap (B \cup C) = (A \cap B) \cup (A \cap C) &\quad
            A \cup (B \cap C) = (A \cup B) \cap (A \cup C)\\
            A \cap \bigcup_{\lambda \in J} B_\lambda = \bigcup_{\lambda
              \in J} (A \cap B_\lambda) &\quad
            A \cup \bigcap_{\lambda \in J} B_\lambda = \bigcap_{\lambda
              \in J} (A \cup B_\lambda) 
        \end{align*}
    \end{proposition}
    Prove some using the definitions
\end{frame}

\begin{frame}{Difference and Complement of Sets}
    \begin{definition}
        \begin{itemize}
            \item The \textbf{difference} of $A$ and $B$ is defined by
            \begin{equation*}
                A \setminus B \;\big(\text{or } A - B) := \{x: x \in A \text{
                  and } x \notin B\}
            \end{equation*}
            \item If $U$ is a set that contains all objects being considered
            in a certain context (a universal set), then $U \setminus A$ is
            called the \textbf{complement} of $A$ in $U$, and is denoted by
            $A^c$
        \end{itemize}        
    \end{definition}
    Examples:
    \begin{itemize}
        \item If $U = \RR$ and $A = (0, 1)$, then $A^c = (-\infty, 0] \cup
        [1, +\infty)$
        \item If $U = [0, 1]$ and $A = (0, 1)$, then $A^c = \{0, 1\}$
    \end{itemize}
\end{frame}

\begin{frame}{Difference and Complement of Sets}
    \begin{proposition}[De Morgan's Laws]
        Let $A$, $B$, and $A_\lambda$ (for $\lambda \in J$) be sets. Then
        \begin{align*}
            (A \cup B)^c = A^c \cap B^c &\qquad (A \cap B)^c = A^c \cup B^c\\
            \left(\bigcup_{\lambda \in J} A_\lambda \right)^c =
            \bigcap_{\lambda \in J} A_\lambda^c &\qquad \left(\bigcap_{\lambda
                  \in J} A_\lambda \right)^c = \bigcup_{\lambda \in J}
            A_\lambda^c
        \end{align*}
    \end{proposition}
    Demonstrate using Venn diagrams
\end{frame}

\begin{frame}[allowframebreaks]{Cartesian Product}
    \begin{definition}
        \begin{itemize}
            \item An \textbf{ordered pair} whose first member is $x$ and
            second member is $y$ is denoted by $(x, y)$
            \item The basic property of ordered pairs is that $(x, y) = (a,
            b)$ iff $x = a$ and $y = b$
            \item The Cartesian product of two sets $A$ and $B$ is defined by
            \begin{equation*}
                A \times B := \{(a, b): a\in A,\, b\in B\}
            \end{equation*}
            \item We can also define $n$-tuples $(a_1, a_2, \ldots, a_n)$.
            Then the Cartesian product of $n$ sets is defined by
            \begin{equation*}
                \prod_{i=1}^n A_i := \{(a_1, a_2, \ldots, a_n): a_i \in A_i,
                \forall i = 1, 2, \ldots, n\}
            \end{equation*}
        \end{itemize}
    \end{definition}
\end{frame}

\begin{frame}
    \frametitle{Cartesian Product}
    \framesubtitle{Examples}
    \begin{itemize}
        \item $\{1, 2\} \times \{3, 4\} = \{(1, 3), (1, 4), (2, 3), (2, 4)\}$
        \item $\RR^n = \underbrace{\RR \times \RR \times \ldots \times
          \RR}_{n \text{ times}}$ 
    \end{itemize}
\end{frame}

\begin{frame}{Common Number Sets}
    \begin{itemize}
        \item Natural numbers: $\NN = \{1, 2, \ldots\}$
        \item Integers: $\ZZ = \{\ldots, -1, 0, 1, \ldots\}$
        \item Rational numbers: $\QQ = \{p/q: p, q \in \ZZ,\, q \neq 0\}$
        \item Real numbers: $\RR$
        \item Complex numbers: $\CC$
    \end{itemize}
\end{frame}



\section{Functions}

\begin{frame}{What is a Function?}
    \begin{definition}
        Let $X$ and $Y$ be two sets. A \textbf{function} or \textbf{map} or
        \textbf{mapping} $f$ from $X$ to $Y$ associates to every element $x
        \in X$ an element $y \in Y$, denoted by $f(x)$. We write
        \begin{align*}
            f\colon X &\to Y\\
            x &\mapsto f(x)
        \end{align*}
        We call $X$ the \textbf{domain} of $f$ and $Y$ the \textbf{codomain} of
        $f$.

        We call $x$ an \textbf{argument} of $f(x)$, and $f(x)$ the \textbf{image}
        of $f$ at $x$.        
    \end{definition}
\end{frame}

\begin{frame}{Image and Inverse Image}
    \begin{definition}
        \begin{itemize}
            \item Let $S \subset X$. The \textbf{image} of $S$ under $f$ is
            defined by $f(S) = \{f(x): x\in S\}$.
            \item The image of $X$ is called the \textbf{range} of $f$.
            \item Let $T \subset Y$. The \textbf{inverse image} of $T$ under
            $f$ is defined by $f^{-1}(T) := \{x: f(x) \in T\}$.
            \item We say $f$ is \textbf{one-one} or \textbf{injective} if for
            every $y \in Y$, there is at most one $x \in X$ such that $y =
            f(x)$.
            \item We say $f$ is \textbf{onto} or \textbf{surjective} if for
            every $y \in Y$, there exists $x \in X$ such that $y = f(x)$.
            \item If $f$ is both one-one and onto, we say $f$ is
            \textbf{bijective}.
        \end{itemize}        
    \end{definition}
\end{frame}

\begin{frame}{Image and Inverse Image}
    \framesubtitle{Examples}
    
\end{frame}

\begin{frame}{Image and Inverse Image}
    \begin{proposition}
        \vspace{-4ex}
        \begin{align*}
            f(C \cup D) = f(C) \cup f(D) &\qquad f\left(\bigcup_{\lambda \in
                  J}A_\lambda\right) = \bigcup_{\lambda \in J} f(A_\lambda)\\
            f(C \cap D) \subset f(C) \cap f(D) &\qquad
            f\left(\bigcap_{\lambda \in J}A_\lambda\right) \subset
            \bigcap_{\lambda \in J} f(A_\lambda)\\
            f^{-1}(C \cup D) = f^{-1}(C) \cup f^{-1}(D) &\qquad
            f^{-1}\left(\bigcup_{\lambda \in J}A_\lambda\right) =
            \bigcup_{\lambda \in J} f^{-1}(A_\lambda)\\
            f^{-1}(C \cap D) = f^{-1}(C) \cap f^{-1}(D) &\qquad
            f^{-1}\left(\bigcap_{\lambda \in J}A_\lambda\right) =
            \bigcap_{\lambda \in J} f^{-1}(A_\lambda)\\
            (f^{-1}(A))^c = f^{-1}(A^c) &\qquad \text{no general results for
              $f$}\\
            A \subset f^{-1}\left(f(A)\right)
        \end{align*}
    \end{proposition}
\end{frame}

\begin{frame}[allowframebreaks]{Composition}
    \begin{definition}
        \begin{itemize}
            \item If $f: X \to Y$ and $g: Y\to Z$, then the \textbf{composition}
            function $g\circ f: X \to Z$ is defined by
            \begin{equation*}
                (g\circ f)(x) = g(f(x))
            \end{equation*}
            \item The \textbf{identity} function on $X$ is defined by
            $Id_X(x) = x$ for all $x \in X$.
            \item We say $f: X \to Y$ is \textbf{invertible} if there exists
            $g: Y \to X$ such that $f \circ g = Id_Y$ and $g \circ f = Id_X$.
            In this case, $g$ is the \textbf{inverse} of $f$ and is denoted
            by $f^{-1}$.
        \end{itemize}        
    \end{definition}
    \begin{itemize}
        \item Example: $f(x) = x^2$, $g(x) = \sin(x)$. Then $(g\circ f)(x) =
        \sin(x^2)$ and $(f\circ g)(x) = \sin^2 (x)$.
        \item The composition operator is associative: $h\circ(g\circ f) =
        (h\circ g)\circ f$.
    \end{itemize}
    \framebreak
    \begin{proposition}
        $f$ is bijective iff $f$ is invertible.
    \end{proposition}
    \emph{Proof:} Suppose that $f$ is bijective. Then, for every $y \in Y$,
    there is a unique $x \in X$ such that $f(x) = y$. This allows us to
    defined $g: Y \to X$ by $f(g(y)) = y$. Therefore, $f\circ g = Id_Y$. Now
    pick any $x \in X$ and let $y = f(x)$. The definition of $g$ implies that
    $g(y) = x$. Therefore, $g(f(x)) = x$ and thus $g \circ f = Id_X$.

    Suppose $f$ is invertible. Then for any $y \in Y$ there is an $x \in X$
    given by $x = f^{-1}(y)$ such that $f(x) = f(f^{-1}(y)) = Id_Y(y) = y$.
    Now suppose there is another $x' \in X$ such that $f(x') = y$. Then
    $f^{-1}(f(x')) = Id_X(x') = x'$. Since $f^{-1}(f(x')) = f^{-1}(y) = x$,
    $x = x'$. Therefore, $f$ is bijective.
\end{frame}

\begin{frame}[allowframebreaks]{Cardinality}
    \begin{definition}
        \begin{itemize}
            \item Two sets $X$ and $Y$ are \textbf{equivalent} if there
            exists a bijection $f: X \to Y$. We write $X \sim Y$.
            \item With every set $A$ we associate a \textbf{cardinal number}
            or \textbf{cardinality} to indicate the ``number'' of elements in
            $A$. Two sets have the same cardinality if they are equivalent.
            \item $Y$ has higher cardinality than $X$ if there exists an
            injection from $X$ into $Y$.
        \end{itemize}
    \end{definition}
    \framebreak
    \begin{definition}
        \begin{itemize}
            \item A set is finite if it is empty or it is equivalent to the
            set $\{1, 2, \ldots, n\}$ for some $n \in \NN$. In this case, it
            has cardinality $n$. A set is \textbf{infinite} if it is not
            finite.
            \item A set is \textbf{denumerable} (or countably infinite) if it
            is equivalent to $\NN$. A set is \textbf{countable} if it is
            denumerable or finite. If a set is denumerable, we say it has
            cardinality $d$.
            \item A set is \textbf{uncountable} if it is not countable. If a
            set is equivalent to $\RR$, we say it has cardinality $c$.
        \end{itemize}
    \end{definition}
    \framebreak
    \begin{proposition}
        \begin{enumerate}
            \item Any subset of a countable set is countable.
            \item The set of rational numbers $\QQ$ is denumerable.
            \item The set of even numbers is equivalent to $\NN$. The set
            $(0, 1)$ is equivalent to $\RR$.
            \item The set $\NN$ and $(0, 1)$ are not equivalent.
            \item The product of two countable sets is countable.
            \item The cardinality of $\RR^n$ is $c$.
            \item The product of two sets of cardinality $c$ has cardinality $c$.
            \item The union of a countable family of countable sets is
            countable.
        \end{enumerate}
    \end{proposition}
\end{frame}

\section{Real Numbers}

\begin{frame}[allowframebreaks]{The Real Number System}
    \begin{itemize}
        \item Most functions we encounter in economics are real-valued
        functions
        \item It is important to know some basic properties of the real
        numbers
        \item It is also key to the study of convergence later on
    \end{itemize}
    \framebreak
    The \textbf{Real Number System} is a set of objects called \textbf{Real
      Numbers} together with two binary operations called \textbf{addition}
    ($+$) and \textbf{multiplication} ($\times$), a binary relation called
    \textbf{less than} ($<$), and two elements called \textbf{zero} (0) and
    \textbf{unity} (1), that satisfies a set of axioms:
    \begin{itemize}
        \item properties of addition
        \begin{enumerate}
            \item [A1] $a + b = b + a$
            \item [A2] $(a + b) + c = a + (b + c)$
            \item [A3] $a + 0 = 0 + a = a$
            \item [A4] there is exactly one real number, denoted by $-a$ such
            that $a + (-a) = (-a) + a = 0$
        \end{enumerate}
        \framebreak
        \item properties of multiplication
        \begin{enumerate}
            \item [A5] $a \times b = b \times a$
            \item [A6] $(a \times b) \times c = a \times (b \times c)$
            \item [A7] $a \times 1 = 1 \times a = a$, and $1 \neq 0$
            \item [A8] if $a \neq 0$ then there is exactly one real number,
            denoted by $a^{-1}$ such that $a \times a^{-1} = a^{-1} \times a = 1$
        \end{enumerate}
        \item the distributive property
        \begin{enumerate}
            \item [A9] $a(b+c) = ab + ac$
        \end{enumerate}
        \item Axioms A1 - A9 are called algebraic axioms
    \end{itemize}
\end{frame}

\begin{frame}[allowframebreaks]
    \frametitle{Consequences of Algebraic Axioms}
    \begin{definition}
        \begin{itemize}
            \item Define \textbf{subtraction} by $a - b = a + (-b)$
            \item Define \textbf{division} by $a/b = a \times b^{-1}$ for $b
            \neq 0$
        \end{itemize}
        
    \end{definition}
    \begin{theorem}
        \begin{enumerate}
            \item If $a + c = b + c$, then $a = b$.
            \item If $ac = bc$ and $c \neq 0$, then $a = b$.
        \end{enumerate}
    \end{theorem}
    \emph{Proof:} Since $a + c = b + c$, $(a + c) + (-c) = (b + c) + (-c)$.
    It follows from A2 that $a + (c + (-c)) = b + (c + (-c))$. By A4, $a + 0
    = b + 0$. By A3, $a = b$. The second claim can be proved in a similar
    way.

    \framebreak
    \begin{theorem}
        If $c\neq 0$ and $d \neq 0$, then
        \begin{enumerate}
            \item $a 0 = 0$, $- (-a) = a$
            \item $(c^{-1})^{-1} = c$, $(-1) a = -a$
            \item $a(-b) = -(ab) = (-a)b$
            \item $(-a) + (-b) = -(a + b)$
            \item $(-a) (-b) = ab$
            \item $(a/c) (b/d) = (ab)/(cd)$
            \item $(a/c) + (b/d) = (ad + bc)/(cd)$
        \end{enumerate}
    \end{theorem}
    Prove some using Axioms A1-A9
\end{frame}

\begin{frame}[allowframebreaks]
    \frametitle{The Real Number System}
    \begin{itemize}
        \item The Order Axioms: there is a subset $P$ of the real numbers,
        called the set of \textbf{positive numbers} such that
        \begin{enumerate}
            \item [A10] For any real number $a$, exactly one of the following
            holds:
            \begin{equation*}
                a = 0 \quad\text{or}\quad a \in P \quad\text{or}\quad -a \in
                P
            \end{equation*}
            \item [A11] If $a \in P$ and $b \in P$, then $a + b \in P$ and
            $ab \in P$
        \end{enumerate}
        \item The ``Less Than'' Relation: define $a < b$ to mean $b - a \in
        P$, $a \leq b$ to mean $b - a \in P$ or $b - a = 0$, $a > b$ to mean
        $a - b \in P$, and $a \geq b$ to mean $a - b \in P$ or $a = b$.
        \item $a < b$ iff $b > a$ and $a \leq b$ iff $b \geq a$
    \end{itemize}
    \framebreak
    \begin{theorem}
        \begin{enumerate}
            \item $a < b$ and $b < c$ implies $a < c$
            \item exactly one of $a < b$, $a = b$, and $a > b$ is true
            \item $a < b$ implies $a + c < b + c$
            \item $a < b$ and $c > 0$ implies $ac < bc$
            \item $a < b$ and $c < 0$ implies $ac > bc$
            \item $0 < 1$ and $-1 < 0$
            \item $a > 0$ implies $1/a > 0$
            \item $0 < a < b$ implies $0 < 1/b < 1/a$
        \end{enumerate}
        Similar results hold for $\leq$
    \end{theorem}
    Prove some using Axioms A1-A11
\end{frame}

\begin{frame}[allowframebreaks]{Completeness Axiom}
    \begin{definition}
        If $S$ is a set of real numbers, then
        \begin{enumerate}
            \item $a$ is an \textbf{upper bound} for $S$ if $x \leq a$ for
            all $x \in S$;
            \item $b$ is the \textbf{least upper bound} (or \textbf{l.u.b.} or
            \textbf{supremum} or sup) for $S$ if $b$ is an upper bound and $b
            \leq a$ whenever $a$ is an upper bound.
            \item We call $b$ a \textbf{maximum} of $S$ if $b$ is the l.u.b.
            for $S$ and $b \in S$.
        \end{enumerate}
        We can similarly define \textbf{lower bound}, \textbf{greatest lower
          bound} (or \textbf{g.l.b.} or \textbf{infimum} or inf), and
        \textbf{mimimum}.

        A set is \textbf{bounded above} if it has an upper bound and is
        \textbf{bounded below} if it has a lower bound.
    \end{definition}
    \framebreak
    \begin{block}{A12 Least Upper Bound Completeness Axiom}
        Suppose $S$ is a nonempty set of real numbers that is bounded above.
        Then $S$ has a l.u.b in $\RR$.
    \end{block}
    \begin{corollary}
        Suppose $S$ is a nonempty set of real numbers that is bounded below.
        Then $S$ has a g.l.b in $\RR$.
    \end{corollary}
    \framebreak
    \begin{proposition}
        Suppose $S$ is a nonempty set of real numbers. Then $b$ is the
        supremum of $S$ iff
        \begin{enumerate}
            \item $x \leq b$ for all $x \in S$, and
            \item for each $\epsilon > 0$, there exists $x\in S$ such that $x
            > b - \epsilon$.
        \end{enumerate}
    \end{proposition}
    \emph{Proof:}
\end{frame}

\begin{frame}{The Archimedean Property}
    \begin{theorem}
        The set $\NN$ is not bounded above.
    \end{theorem}
    \emph{Proof:} Suppose $\NN$ is bounded above. By the Completeness Axiom,
    there is a supremum $b \in \RR$ such that $n \leq b$ for all $n \in \NN$.
    Since $n+1 \in \NN$, $n + 1 \leq b$ and thus $n \leq b - 1$ for all $n
    \in \NN$. This is a contradiction since $b$ is no longer the least upper
    bound of $\NN$.
    \vspace{2ex}

    \begin{theorem}
        \begin{enumerate}
            \item For any $x, y \in \RR$, if $x < y$ then there exists a
            rational (irrational) number $r$ such that $x < r < y$.
            \item For any $x \in \RR$ and $\epsilon > 0$, there exists a
            rational (irrational) number $r$ such that $0 < |r - x| <
            \epsilon$.
        \end{enumerate}
    \end{theorem}
\end{frame}

\begin{frame}{The Archimedean Property}
    \framesubtitle{Exercise}
    Prove that the supremum of $\{x \in \RR: x < 3\}$ is 3.

    \emph{Proof:}
\end{frame}

\begin{frame}[allowframebreaks]{Boundedness and Monotonicity of Functions}
    \begin{definition}
        Let $f: X \to Y$ be a \textbf{real-valued} function, that is, $Y
        \subset \RR$.

        Define upper bounds, the supremum, the maximum of $f$ as the upper
        bounds, supremum, and maximum of its range $f(X) \subset \RR$. The
        arguments $x$ such that $f(x)$ is the maximum are called the
        \textbf{maximizers} of $f$.

        We say $f$ is \textbf{bounded} if its range is bounded.
    \end{definition}
    \framebreak
    \begin{definition}
        If $X \subset \RR$, then we say
        \begin{itemize}
            \item $f$ is \textbf{increasing} if $x \leq x' \implies f(x) \leq
            f(x')$
            \item $f$ is \textbf{decreasing} if $x \leq x' \implies f(x) \geq
            f(x')$
            \item $f$ is \textbf{monotonic} if it is increasing or decreasing
            \item $f$ is \textbf{strictly increasing} if $x < x' \implies f(x) <
            f(x')$
            \item $f$ is \textbf{strictly decreasing} if $x < x' \implies f(x)
            > f(x')$
        \end{itemize}
    \end{definition}
\end{frame}


\end{document}